% PROPOSAL (P-09) — <=500-word version, SiBi/UFPR norm.
% Does NOT replace 0-iniciais/abstract.tex. The cut is the author's editorial call.
%
% WHAT THIS VERSION IS, measured rather than claimed: a FAITHFUL CONDENSATION,
% not an excerpt. Of its 19 sentences, 3 are verbatim from the full abstract and
% 16 were rewritten — sentences merged and clauses compressed to fit the cap. An
% earlier version of this header said "removed, never rewritten"; that was FALSE,
% and the measurement by this proposal's own author refuted it.
%
% What did NOT change, and this remains script-proven: no new number (the short
% version's number set is a SUBSET of the full one), no claim absent from the
% full version, and the cut is identical in PT and EN.
\begin{abstract}

Short-text classification in Portuguese depends on large volumes of labels, and that labeling is the dominant cost of these systems. The case studied is that of product descriptions from electronic invoices, with abbreviated vocabulary and hundreds of imbalanced categories. This thesis proposes and evaluates FALCO (Active Learning Framework with LLMs for Short Text), which attacks that cost on three fronts: informed composition of the initial set, uncertainty-based active selection, and partial replacement of the human annotator by oracles based on large language models (LLMs), in phases of increasing cost. The investigation uses an audited public dataset of 250,365 descriptions across 621 categories. The protocol is pre-registered, with immutable partitions, paired analysis (McNemar, Wilcoxon), and a traceable artifact for every number. Five results support the thesis. First, the composition of the initial set matters more the smaller the budget, and the proposed heuristic, DRI-SL, outperforms without using any label the best individual found by a supervised genetic algorithm at every size from 100 to 5,000. Second, zero-shot LLM oracles reach $77$--$83\%$ accuracy in the closed space of 621 categories, at costs of US\$~0.035 to US\$~0.92 per thousand labels, with Macro F1 above that of a light supervised classifier trained on the full base. Third, the measurement instrument is part of the result: without output restriction, $6.8\%$ of responses become spurious errors, and the same free model diverges significantly from itself across two serving providers ($p<0.001$). Fourth, uncertainty strategies outperform random selection at the maximum significance achievable with 8 seeds ($p=0.0078$), and the advantage survives uniform oracle noise up to $\varepsilon=0.4$, a range that covers the error of the real LLMs. No oracle reached the pre-registered $85\%$ accuracy criterion, which made learning with noisy labels the central scenario. Fifth, the framework was executed end to end with a free LLM oracle, at zero monetary cost, and validated with BERTimbau across three seeds. The central hypothesis was to obtain at least $95\%$ of the accuracy of full supervision on the reference pool with at most 34,724 labels. The verdict is twofold. \textbf{The criterion is attainable within the ceiling, and is crossed at 20,000 labels ($8.6\%$ of the population) in all three seeds, in a \textit{post hoc} gold-label sweep. The executed configuration, however, stopped by stagnation short of the floor, at 11,936 labels ($5.2\%$)}. The paired decomposition exonerates oracle noise and selection, and attributes the loss to the stopping criterion inherited from the light classifier. The limitations are declared: the results come from a single dataset, from the author's own study domain, and the cost analysis rests on ratios between oracles, not on absolute prices. The thesis also delivers the LCE (Learning Curve Efficiency) metric, the open-source \texttt{activelearning} library, and the FlowBuilder tool, which make the protocol reusable.

\end{abstract}
